\documentclass[11pt]{article}

\usepackage{amsmath,amssymb}
\usepackage[margin=1in]{geometry}
\usepackage{hyperref}

\title{Product Returns under Perfect Competition}
\author{Draft notes}
\date{\today}

\begin{document}
\maketitle

This note records a competitive private-ownership economy with a continuum
of buyers and a continuum of identical price-taking firms. The first section
is the environment without returns. The second allows inspection and a refund.
Unlike the incidence notes, the refund is not restricted to equal the list
price.

\section{No returns}

\subsection{Setting}

A continuum of consumers and a continuum of identical price-taking firms
trade one good that cannot be returned. Each firm produces at constant unit
cost $c$. Consumers have unit demand and quasi-linear preferences, with an
outside option of zero, and differ in the value $v_i$ they know when they
buy.

They take a price $p$ as given. Consumer $i$ buys when $v_i\geq p$. A firm
that sells $y$ units earns $(p-c)y$.

\subsection{Equilibrium}

The firm supplies nothing if $p<c$, any scale if $p=c$, and infinite output
if $p>c$. Demand is the measure of consumers with $v_i\geq p$, and that
measure is finite. Positive demand therefore clears only at $p=c$. Profit
is zero, and firms meet that demand.

\section{Returns}

\subsection{Setting}

The good can now be purchased, inspected, and returned. Checkout price and
refund are two prices. Each firm still produces at cost $c$ and now incurs
a unit handling cost $h$ on every returned item. Returned units are salvaged
rather than restocked.

Consumers still have unit demand, quasi-linear risk-neutral preferences, and
an outside option of zero. They differ in return hassle $\kappa_i$, the
disutility of sending the good back, and in the distribution $G_i$ of match
value $v$ realized after purchase. Consumer $i$ knows $(\kappa_i,G_i)$ when
she buys, but not $v$.

Consumers take a checkout price $p$ and a refund $R$ as given. Each buys at
most one unit, or not at all, then draws $v\sim G_i$ and decides whether to
keep the item or return it.

If she has bought at $(p,R)$, keeping yields $v-p$ and returning yields
$R-p-\kappa_i$. She therefore keeps if and only if $v\geq R-\kappa_i$. The
keep rule depends on the refund and on hassle, not on the checkout price.
Expected surplus from buying is
\[
  cs_i(p,R)
  =
  \mathbb{E}_{v\sim G_i}\bigl[\max(v-p,\,R-p-\kappa_i)\bigr]
  =
  -p+\mathbb{E}_{v\sim G_i}\bigl[\max(v,\,R-\kappa_i)\bigr].
\]
She buys when that value is nonnegative. The reservation checkout price for
a given refund, $p_i(R)=\mathbb{E}[\max(v,R-\kappa_i)]$, is the largest $p$
at which she is willing to buy. Heterogeneity therefore cannot be collapsed
into a single willingness to pay that is independent of the refund. Consumer
$i$'s return probability at refund $R$ is $\phi_i(R)=G_i(R-\kappa_i)$.

A firm that sells $y$ units collects $p$ on each sale, refunds $R$ and
incurs $h$ on the share $r$ that come back, and therefore earns
\[
  \pi=\bigl(\mu-c\bigr)y,
  \qquad
  \mu=p-r(R+h),
\]
where $\mu$ is the effective margin. A price-taker then compares $\mu$ with
$c$.

Industry output equals the measure of purchases. The returned share $r$ is
the mean of $\phi_i(R)$ among those who buy, and who buys depends on $(p,R)$.

\subsection{Equilibrium}

A competitive equilibrium is a pair $(p,R)$, purchases, a return rate $r$,
and industry output $y$ such that consumers optimize given $(p,R)$,
firms maximize profit given $(p,R)$ and $r$, and the market clears.
Both sides see the same $(p,R)$.

Consumer $i$ buys if and only if $cs_i(p,R)\geq 0$, and returns with
probability $\phi_i(R)$. The return rate $r(p,R)$ is the mean of
$\phi_i(R)$ among those buyers. A firm takes $(p,R)$ and $r$ as given and
earns $(\mu-c)y$, so it supplies nothing if $\mu<c$, any scale if $\mu=c$,
and infinite output if $\mu>c$.

Market clearing requires $y$ to equal the measure of buyers, and the $r$ in
the effective margin must be $r(p,R)$. Infinite supply cannot clear against
a finite mass of consumers, and zero supply cannot clear if a positive
measure wants to buy. Whenever demand is positive, therefore
\begin{equation}
  \label{eq:returns_ce}
  p-r(p,R)\,(R+h)=c.
\end{equation}
As long as this holds, the pair is a competitive equilibrium. Consumers take
$(p,R)$ as given and follow the buy and keep rules, which determines demand
and $r(p,R)$. Then $\mu=c$, so every firm is willing to sell any scale, and
industry can meet that demand. Buyers optimize, sellers optimize, and the
market clears.

\eqref{eq:returns_ce} is one equation in two prices, so many pairs are
competitive equilibria. The buyer's keep rule agrees with the social rule,
keep when $v\geq -h-\kappa_i$, only at $R=-h$. Zero profit then requires
$p=c$. That pair satisfies \eqref{eq:returns_ce}, so the efficient
allocation can be sustained as a competitive equilibrium. At other pairs
satisfying \eqref{eq:returns_ce}, consumers return too often or too rarely.

The First Welfare Theorem fails because the firm cannot choose the quantity
of returned units. It must take back whatever the buyer returns, and its
only quantity choice is how many new units to supply. Given the checkout
price and the refund, the buyer determines the demand for the new good and
the supply of used goods. She demands the new good when her expected surplus
is nonnegative and supplies a used unit when $v<R-\kappa_i$. The firm takes
the resulting return rate as given and compares the effective margin with
$c$.

If the firm could choose how many returned units to buy, it would gain on
each return when $R<-h$ and lose when $R>-h$. The return market would
therefore clear only at $R=-h$, and production would clear only at $\mu=c$.
At that refund the effective margin equals the checkout price, so $\mu=c$
becomes $p=c$. The firm then earns zero on each new unit and zero on each
returned unit, so it is willing to produce any quantity and to take back any
quantity. Any demand for the new good and any supply of used goods can
therefore clear.

Adding a resale market can raise welfare. Used units then trade at a price
the firm can bid for, and clearing selects only the efficient allocation.

\end{document}
